\documentclass[10pt,a4paper]{amsart}
\usepackage[margin=19mm]{geometry}
\usepackage{amsmath,amssymb,mathtools}
\usepackage[scr=rsfs,cal=euler]{mathalfa}
\usepackage{xfrac}
\usepackage[unicode,hidelinks,bookmarks=false]{hyperref}
\newcommand{\ie}{i.e.\ }
\newcommand{\cf}{cf.\ }
\newcommand{\resp}{respectively\ }
\newcommand{\Subsection}{\subsection}
\providecommand{\nobreakdash}{\nobreak\hbox{-}}
\setlength{\emergencystretch}{2em}
\multlinegap0pt
\usepackage{amssymb}
\newtheorem{prop}{Proposition}
\newtheorem{lemm}{Lemma}
\title{Rank $2$ vector bundles and degrees of points of del Pezzo surfaces}
\author{Claire Voisin}
\date{Source: arXiv:2509.17996v2 (CC BY 4.0)}
\begin{document}
\maketitle

\noindent\emph{Source and licence.} Rank $2$ vector bundles and degrees of points of del Pezzo surfaces. Version arXiv:2509.17996v2 (CC BY 4.0), distributed by arXiv under the Creative Commons Attribution 4.0 International licence, http://creativecommons.org/licenses/by/4.0/, as recorded in the arXiv licence metadata for this exact version and shown on the abstract page https://arxiv.org/abs/2509.17996v2. Full licence notice: \texttt{licenses/arxiv-2509.17996-CC-BY-4.0.txt}.\par\smallskip

\noindent\textit{Editorial scope.} Counted unit: the article's prop prodominant (source0.tex lines 203--206). The article's complete original proof of it is delivered with it (source0.tex lines 211--296, one \texttt{proof} environment of 86 lines, which the article places away from the statement). No mathematics is written, repaired or re-derived: the statement and the proof are the article's own lines, in the article's own notation, and no retained line is substituted or re-flowed.  No further source-local context is needed: every cross-reference of every retained line resolves inside the two delivered spans.  Nothing was omitted from any retained span, so the package carries no omission record.  The retained lines cite 1 works, whose entries are transcribed verbatim from the article's own bibliography source in the e-print and delivered with short editorial labels recorded in the item's reference mapping.  No retained line contains a figure environment, an image-inclusion command or drawing code, so no image is delivered and the package declares no asset dependency.  Editorial formatting only, in addition: an AMS \texttt{amsart} preamble that restates the article's own declarations and macros used by the retained lines, namely the environments \texttt{prop}, \texttt{lemm} on separate counters, together with the standard packages those lines need.  The retained environments print in this document as Proposition 1, Lemm 1 in order of appearance, whereas the article prints the same items with the numbering of its own full document.  The complete native source of the article as distributed in the arXiv e-print for this exact version is bundled unchanged as \texttt{sources/185501/source0.tex}.
\begin{prop}\label{prodominant}  Assume $n\geq 3$ and ${\rm char}\,K=0$. Then, for  $s=-1$, the rational map
$$\Psi_s: G\times G\dashrightarrow X^{[3]}$$
is dominant.
\end{prop}

\begin{proof}[Proof of Proposition \ref{prodominant}] For $s=-1$, the morphism
$$\Phi_s: G\times X\dashrightarrow X,$$
$$\Phi_s([W],x)=\mu_{s,W}(x)$$
has a special form, namely,  on each hyperplane section $E\subset X$, the map $\mu_{E,-1}$ is the multiplication by $-2$ on the elliptic curve $E$ and  maps $x\in E$ to
$h_E-2x$, where $h_E:=c_1(\mathcal{O}_E(1))$, hence is geometrically realized by sending $x$ to the residual intersection point of the projective  line
$\mathbb{P}(T_{E,x})$ tangent to $E$ at $x$ with $X$. It follows that $\mu_{s,W}(x)$ depends only on the tangent space at $x$ of the curve $E_{W,x}$ passing through $x$. In other words, the rational map $\Phi_{-1}$ factors
  through the rational morphism
$$ G\times X\dashrightarrow\mathbb{P}(T_X)$$
which to $([W], x)$ associates the tangent line to the fiber $E_{W,x}$ passing through $x$  of the linear projection $\phi_W$.

We observe that it suffices to prove the result when  ${\rm dim}\,X=2$ since any set of three points on $X$ (or rather subscheme of length $3$) is supported on a smooth cubic surface.
Furthermore we can assume that the field is algebraically closed (for example $K=\mathbb{C}$). We choose three  general points $x,\,y,\,z$ on $X$ and have to prove that the preimage
$\Psi_{-1}^{-1}(\{x,\,y,\,z\})$ is not empty.
Looking at the construction of $\Psi_{-1}$, this preimage consists of a pencil $W$ of elliptic plane curves on $S$, and a set of three {\it collinear} points $x',\,y',\,z'$ on $X$ (generating a line $\Delta_{W'}$), such that, denoting respectively
$E_x,\,E_y,\,E_z$ the fibers of the pencil passing through $x,\,y,\,z$, we have

\begin{eqnarray}\label{eqpointcourbe} x'\in E_x,\,y'\in E_y,\,z'\in E_z,\\
\label{eqpointmultmoinsun}\mu_{E_x,-1}(x')=x,\, \mu_{E_y,-1}(y')=y,\,\mu_{E_z,-1}(z')=z.\end{eqnarray}


Using the above description of the maps $\mu_{E,-1}$, we can describe differently this fiber, namely, let
$C_x\subset X$ (resp. $C_y\subset X$, resp. $C_z\subset X$) be the curve of points $x'\in X$ (resp. $y'\in X$, resp. $z'\in X$)  such that the line
$\langle x',\,x\rangle$ (resp. $\langle y',\,y\rangle$, resp. $\langle z',\,z\rangle$) is tangent to $X$ at $x'$  resp. $y'$, resp. $z'$ (more rigorously, we should take the respective Zariski closures of these curves in $X\setminus \{x\}$,  $X\setminus \{y\}$ and $X\setminus \{z\}$). These curves, which are ramification curves of the projection of $X$ to $\mathbb{P}^2$  from $x$ (resp. $y$, resp. $z$), are well understood. They are members of the linear system  $|\mathcal{O}_X(2)|$, irreducible for general points $x,\,y,\,z$, and they contain respectively the points $x,\,y,\,z$. Furthermore, they are mobile in the sense that  any point $w$ of $X$ can be chosen not to belong to $C_x$ (resp. $C_y$, resp. $C_z$). This last point is clear since there exists a line $\Delta$ passing through $w$ and not tangent to $X$ at $w$. This line intersects $X$ at another point $x\in X$, and $w$ does not belong to $C_x$.

 For $(x',\,y',\,z')\in C_x\times C_y\times C_z$, we then choose  any elliptic plane curve $E_{x',x}$  passing through $x'$ and $x$, and similarly $E_{y',y}$  passing through $y'$ and $y$, $E_{z',z}$  passing through $z'$ and $z$.
Equations (\ref{eqpointcourbe}) and (\ref{eqpointmultmoinsun}) are then automatically satisfied since the line $\langle x',\,x\rangle$ is tangent to any such $E_x$ at $x'$ and similarly for $y$ and $z$. We need now  to impose the following  conditions~:
\begin{enumerate}
\item \label{it1} The three points $x',\,y',\,z'$ are collinear (producing the line $\Delta_{W'}$).
\item \label{it2} The three elliptic curves $E_{x',x}$, $E_{y',y}$, $E_{z',z}$ generate a pencil (producing the  pencil $W$).

\end{enumerate}
Condition \ref{it1} has to be satisfied on $C_x\times C_y\times C_z$. Given $x',\,y',\,z'$, Condition \ref{it2} has to be satisfied on
the product $\mathbb{P}^1_{x',x}\times \mathbb{P}^1_{y',y}\times \mathbb{P}^1_{z',z}$, where the line $\mathbb{P}^1_{x',x}$  parameterizes planes in $\mathbb{P}^3$ containing the line $\langle x',x\rangle$  and so on.

That the set of triples $(x',\,y',\,z')\in C_x\times C_y\times C_z$ of {\it distinct } points satisfying conditions \ref{it1} and \ref{it2} is not empty  follows now from the following
\begin{lemm} (1) For a general choice of points $x,\,y,\,z$, condition \ref{it1} is satisfied along a nonempty curve $D\subset C_x\times C_y\times C_z$ with the property that, for a general triple $(x',\,y',\,z')\in D$, we have $x\not=x',\,y\not=y',\,z\not=z'$ and the three lines  $\langle x',x\rangle$, $\langle y',y\rangle$, $\langle z',z\rangle$ are mutually non-intersecting.

(2) For a general point $(x',\,y',\,z')\in D$, condition \ref{it2} is satisfied along a curve $D'_{x',\,y',\,z'}\subset \mathbb{P}^1_{x',x}\times \mathbb{P}^1_{y',y}\times \mathbb{P}^1_{z',z}$ which is isomorphic to $\mathbb{P}^1$.
\end{lemm}
\begin{proof} (1) The set $D$ cannot contain a surface. Indeed, it would dominate otherwise one of the 3 products $C_x\times C_y$, $C_x\times C_z$, $C_y\times C_z$ under the corresponding projections. Assume it dominates $C_x\times C_y$. Then for any $(x',\,y')\in C_x\times C_y$, the third intersection point $z''$ of the line $\langle x',\,y'\rangle$ with $X$ must belong to the curve $C_z$. Using the mobility of  the curve $C_z$  with $z$ as explained above, this is not possible if $z$ is chosen generically. Next, the set $D$ has expected codimension $\leq 2$ in $C_x\times C_y\times C_z$. Indeed, on each curve $C_\bullet$, where  $\bullet=x,\,y,\,z$, we have
the inclusion
$$\mathcal{O}_{C_\bullet} (-1)\hookrightarrow W_4\otimes \mathcal{O}_{C_\bullet},$$
where $W_4=H^0(X,\mathcal{O}_X(1))^*$.
We combine these three inclusion maps to construct a morphism of vector bundles
\begin{eqnarray}\label{eqmophvnun}\phi: {\rm pr}_x^*\mathcal{O}_{C_x} (-1)\oplus {\rm pr}_y^*\mathcal{O}_{C_y} (-1)\oplus {\rm pr}_z^*\mathcal{O}_{C_z} (-1)\rightarrow W_4\otimes \mathcal{O}_{C_x\times C_y\times C_z},\end{eqnarray}
of vector bundles of respective ranks $3$ and $4$.
It is clear that the locus $D$ of collinear triples $(x',\,y',\,z')$ is contained in the locus $D_2$ where $\phi$ has rank $\leq 2$, but  we have to remove from it the sublocus $D'_2$ where $x'=y'$ or $x'=z'$, or $y'=z'$. It is a standard fact that the rank locus $D_2$ has expected codimension $\leq 2$, hence its codimension is exactly $2$ by the previous assertion. The class of $D_2$ in ${\rm CH}^2(C_x\times C_y\times C_z)$ is in fact computed following \cite{fulton}. This class is nothing but the Segre class
$s_2(\mathcal{F})$ in ${\rm CH}^2(C_x\times C_y\times C_z)$, where $\mathcal{F}:={\rm pr}_x^*\mathcal{O}_{C_x} (-1)\oplus {\rm pr}_y^*\mathcal{O}_{C_y} (-1)\oplus {\rm pr}_z^*\mathcal{O}_{C_z} (-1)$. This follows indeed from  the fact that the rank $\leq 2$ locus of $\phi$ is also the image in $C_x\times C_y\times C_z$, under the projection map $\pi: \mathbb{P}(\mathcal{F})\rightarrow C_x\times C_y\times C_z$, of the locus
$\widetilde{D}_2\subset \mathbb{P}(\mathcal{F})$ defined by the $4$ sections of $\mathcal{F}^*$ or sections of $\mathcal{O}_{\mathbb{P}(\mathcal{F})}(1)$ given by $ \phi$. Using the fact that $D_2$ has the right dimension, this is saying that
$$[D_2]=\pi_* (c_1(\mathcal{O}_{\mathbb{P}(\mathcal{F})}(1))^4)=s_2(\mathcal{F})\,\,{\rm in}\,\,{\rm CH}^2(C_x\times C_y\times C_z).$$
The degree  of $D_2$ is computed using the Whitney formula for  the total Segre class, which gives
$$ s(\mathcal{F})={\rm pr}_x^*( 1+c_1(\mathcal{O}_{C_x}(1))){\rm pr}_y^*( 1+c_1(\mathcal{O}_{C_y}(1))){\rm pr}_z^*( 1+c_1(\mathcal{O}_{C_z}(1))) \,\,{\rm in}\,\,{\rm CH}(C_x\times C_y\times C_z),$$
so that
\begin{eqnarray} \label{eqpoursegre2} s_2(\mathcal{F})={\rm pr}_x^*(c_1(\mathcal{O}_{C_x}(1))){\rm pr}_y^*( c_1(\mathcal{O}_{C_y}(1)))+{\rm pr}_x^*(c_1(\mathcal{O}_{C_x}(1))){\rm pr}_z^*(c_1(\mathcal{O}_{C_z}(1)))\\
\nonumber +{\rm pr}_y^*(c_1(\mathcal{O}_{C_y}(1))){\rm pr}_z^*(c_1(\mathcal{O}_{C_z}(1))) \,\,{\rm in}\,\,{\rm CH}(C_x\times C_y\times C_z).
\end{eqnarray}
In order  to prove the non-emptiness of $D_2\setminus D'_2$, it suffices to  show that
$${\rm deg}_{H_x}\,D_2> {\rm deg}_{H_x}\,D'_2,$$ where the degree is computed via the line bundle $H_x:={\rm pr}_x^*\mathcal{O}_{C_x}(1))$ on $C_x\times C_y\times C_z$.
The curves $C_x,\,C_y,\,C_z$ being defined by quadratic equations in $X$, we have  $${\rm deg}_{H_x}\,D'_2={\rm deg}(C_y\cdot C_z){\rm deg}_{H_x}(C_x)=12\cdot 6,$$
$${\rm deg}_{H_x}\,D_2=6\cdot 6\cdot 6,$$
from which we conclude that ${\rm deg}_{H_x}\,D_2> {\rm deg}_{H_x}\,D'_2$ and $D_2\setminus D'_2$ is non-empty.

We prove by a similar counting argument that for a general element $(x',\,y',\,z')$ of $D_2\setminus D'_2$, the three lines \begin{eqnarray}\label{eq3lines}\langle x',x\rangle, \,\langle y',y\rangle, \,\langle z',z\rangle \end{eqnarray} are mutually non-intersecting.




 (2) will  follow from the last statement. Indeed, this statement says that the  three points $\delta_{x'x},\,\delta_{y'y},\,\delta_{z'z}$ which parameterize respectively the three  lines
(ref{eq3lines})
are three points in general position in the Grassmannian of lines in $\mathbb{P}^3$. The Grassmannian $G(2,4)$ is a quadric $Q$ in $\mathbb{P}^5$, and two lines in $\mathbb{P}^3$ are non-intersecting if and only if the corresponding points in $G(2,4)=Q$ have the property that the line they generate is not contained in $Q$.   There is  thus a single orbit  under the action of ${\rm PGL}(4)$  on $G(2,4)^{[3]}$  of triples parameterizing three lines  mutually nonintersecting, as it follows from the similar statement for the action of the orthogonal group ${\rm O}(6)$ on $Q$. It thus suffices to prove the result when the three lines (\ref{eq3lines}) in $\mathbb{P}^3$  are defined  by equations
$$X_0=X_1=0,\,X_2=X_3=0,\,X_0-X_2=0,\,X_1-X_3=0.$$
An easy computation shows that the set of triples of planes
\begin{eqnarray}\label{eq3planes} p_{x',x}, \,p_{y',y}, \,p_{ z',z}\in (\mathbb{P}^3)^*
\end{eqnarray}
such that the corresponding plane $P_{x',x}$ contains $\langle x',\,x\rangle$ and so on, and such that the three linear forms $p_{x',x}, \,p_{y',y}, \,p_{ z',z}$ generate a pencil of planes (i.e. are collinear),  is a copy of $\mathbb{P}^1$ diagonally embedded in $\mathbb{P}^1_{x',x}\times \mathbb{P}^1_{y',y}\times \mathbb{P}^1_{z',z}$. Indeed,
one writes   $$p_{x',x}=u_0X_0+u_1 X_1,\,p_{y,y'}=v_2X_2+v_3X_3,$$
 $$p_{z,z'}=w(X_0-X_2)+w'(X_1-X_3),$$ for adequate choice of homogeneous coordinates on $\mathbb{P}^3$. The fact that the three planes generate a pencil then provides  equations
$$ w=au_0=-bv_2,\,w'=au_1=-bv_3
$$
for some nonzero coefficients $a,\,b$. Hence the equations provide
$$u_1/v_3=u_0/v_2=-b/a$$
$$w=au_0,\,w'=au_1,$$
which proves the result.
\end{proof}
This concludes the proof of Proposition \ref{prodominant}.
\end{proof}

\begin{thebibliography}{99}
\bibitem[fulton]{fulton} W. Fulton. Intersection theory. Ergebnisse der Mathematik und ihrer Grenzgebiete (3) [Results in Mathematics and Related Areas (3)], 2. Springer-Verlag, Berlin, 1984.
\end{thebibliography}
\end{document}
